% system-model.tex -- BNR-2026-012

\section{System Model}
\label{sec:system-model}

\subsection{Architecture Overview}

The production DAC operates within the Basis Network zkEVM L2
architecture. Each enterprise operates a dedicated L2 chain with a
centralized sequencer that batches transactions and submits validity
proofs to the Basis Network L1 (an Avalanche Subnet-EVM chain). The DAC
sits between the sequencer and the L1: the sequencer disperses batch
data to the committee before proof submission, and the committee's
attestation is verified on-chain alongside the ZK proof.

\subsection{Committee Configuration}

The committee consists of $n = 7$ enterprise-operated nodes with a
reconstruction and attestation threshold of $k = 5$:

\begin{itemize}
  \item \textbf{Fault tolerance}: up to 2 node failures (crash or
    Byzantine).
  \item \textbf{Honest minority}: 3 honest nodes suffice to block false
    attestation ($n - k + 1 = 3$).
  \item \textbf{Privacy threshold}: fewer than 5 colluding nodes cannot
    reconstruct the AES key (Shamir guarantee).
  \item \textbf{Fallback}: if fewer than 5 nodes are available, the
    system falls back to on-chain data availability (AnyTrust
    model~\cite{anytrust}).
\end{itemize}

\subsection{Threat Model}

We consider the following adversarial capabilities:

\begin{enumerate}
  \item \textbf{Passive corruption}: up to $k - 1 = 4$ committee
    members may collude to attempt data reconstruction. The AES-256-GCM
    encryption and Shamir key sharing ensure that $k - 1$ shares reveal
    zero information about the key.
  \item \textbf{Crash failures}: up to $n - k = 2$ nodes may
    simultaneously become unavailable. Reed-Solomon $(5,7)$ coding
    ensures reconstruction from any 5 surviving nodes.
  \item \textbf{Malicious disperser}: the sequencer may attempt to
    distribute invalid RS chunks. KZG polynomial commitments (planned
    for production) enable each node to verify its chunk before attesting.
  \item \textbf{Network adversary}: an adversary may delay or reorder
    messages between the disperser and committee members. The protocol
    assumes eventual delivery within a configurable timeout.
\end{enumerate}

\subsection{Hybrid Privacy Model}

The protocol provides two layers of privacy:

\begin{itemize}
  \item \textbf{Computational data privacy}: batch data is encrypted
    with AES-256-GCM using a random 254-bit key (reduced modulo the
    BN254 scalar field prime). Breaking data privacy requires breaking
    AES-256-GCM, which is computationally infeasible under standard
    assumptions.
  \item \textbf{Information-theoretic key secrecy}: the AES key is
    shared via Shamir $(5,7)$-SS over the BN254 scalar field. Any
    coalition of fewer than 5 shareholders learns zero information about
    the key, regardless of computational power.
\end{itemize}

\subsection{Invariants}
\label{sec:invariants}

We formalize six invariants that the production DAC must satisfy. These
invariants serve as the specification interface for downstream TLA+
formalization and Coq verification.

\textbf{INV-DA-P1 (Verifiable Encoding).}
Every RS chunk distributed to DAC members must be verifiable against a
polynomial commitment. A node receiving an invalid chunk must reject it
and not attest.

\textbf{INV-DA-P2 (Data Recoverability).}
If at least $k = 5$ of $n = 7$ nodes store valid chunks, the complete
batch data must be recoverable by: (a) collecting any 5 chunks,
(b) RS-decoding to ciphertext, (c) collecting 5 key shares,
(d) SSS-recovering the AES key, (e) decrypting.

\textbf{INV-DA-P3 (Enterprise Privacy).}
No individual DAC node can reconstruct the batch data from its chunk and
key share alone. A node's chunk is encrypted ciphertext (AES-256-GCM),
and a single Shamir share reveals zero information about the AES key
(information-theoretic guarantee).

\textbf{INV-DA-P4 (Attestation Liveness).}
If at least 5 of 7 nodes are online and the batch data is correctly
encoded, attestation must complete within the configured timeout
(default: 1 second).

\textbf{INV-DA-P5 (Fallback Safety).}
If fewer than 5 nodes are available, the system must not attest. Instead,
the system falls back to on-chain DA (posting batch data directly to L1).

\textbf{INV-DA-P6 (Availability Guarantee).}
With per-node availability $p \geq 0.99$, the probability that the batch
data is available (at least $k = 5$ of $n = 7$ nodes online) must exceed
$99.999\%$.

\subsection{Availability Analysis}

For a $(k,n)$ DAC where each node has independent availability
probability $p$, the system availability is:

\begin{equation}
P(\text{available}) = \sum_{i=k}^{n} \binom{n}{i} p^i (1-p)^{n-i}
\label{eq:availability}
\end{equation}

Table~\ref{tab:availability-configs} presents the availability for
different configurations and node reliability levels.

\begin{table}[t]
\centering
\caption{Availability by Configuration and Node Reliability}
\label{tab:availability-configs}
\begin{tabular}{@{}lccccc@{}}
\toprule
\textbf{Config} & $k$ & $n$ & $p{=}0.95$ & $p{=}0.99$ & $p{=}0.999$ \\
\midrule
RU-V6 (2,3)     & 2 & 3 & 99.3\%  & 99.97\%   & 99.9997\% \\
\textbf{Proposed (5,7)} & \textbf{5} & \textbf{7} & \textbf{99.6\%} & \textbf{99.997\%} & \textbf{99.99999\%} \\
AnyTrust (5,6)  & 5 & 6 & 96.7\%  & 99.85\%   & 99.999\%  \\
Conservative (4,7) & 4 & 7 & 99.98\% & 99.99997\% & $\approx$1.0 \\
\midrule
\textbf{Nines} & & & & & \\
RU-V6 (2,3)     & & & 2.1 & 3.5 & 5.5 \\
\textbf{Proposed (5,7)} & & & \textbf{2.4} & \textbf{4.5} & \textbf{7.5} \\
AnyTrust (5,6)  & & & 1.5 & 2.8 & 4.8 \\
Conservative (4,7) & & & 3.7 & 6.5 & 10+ \\
\bottomrule
\end{tabular}
\end{table}

At $p = 0.99$ (enterprise-grade infrastructure), the $(5,7)$
configuration achieves 4.5 nines of availability (99.997\%), exceeding
the 3-nine target (99.9\%) specified in INV-DA-P6. At $p = 0.999$,
availability reaches 7.5 nines (99.99999\%). The $(5,7)$ configuration
strictly dominates the $(2,3)$ baseline at $p \geq 0.95$ despite the
higher threshold, because the additional nodes provide combinatorial
redundancy that outweighs the increased threshold requirement.
